\chapter{The product decision}
\label{ch:purpose}

A company's accounts describe the financial consequences of its activities,
but they leave much of the activity itself unobserved. Revenue combines many
sales, a borrowing balance combines contracts with different terms, and a
provision combines estimates about uncertain future payments. Aggregation is
necessary for reporting. It also makes the inverse task---recovering economic
history from published reports---fundamentally different from entering known
transactions into a ledger.

AccDSL will preserve this distinction. Its first product is an evidence-backed
representation of a company's historical record. Its second is a set of
compatible histories and current economic states. Its third is a collection of
conditional simulations, each explaining how a particular history, accounting
treatment and business response produces future journals and reports. These
products share identifiers and provenance, so an analyst can move from a
forecast balance to the assumption or disclosure responsible for it.

The recommended design is a small, typed accounting language over a shared
intermediate representation, combined with an evidence database, contract
calculators and replaceable inference engines. A language earns its place by
making recurring accounting decisions explicit and checkable: which entity and
period an amount belongs to, whether it is known or assumed, which rule applies,
and how an event changes a balanced ledger. The language does not supply missing
facts or determine a company's future behaviour.

\section{Why journals are necessary but insufficient}

Consider a company with a borrowing of 100 currency units and annual interest
expense of 5. Assume that the entire principal was outstanding throughout the
year, interest was paid at year-end, and there were no fees, defaults, hedges or
foreign-exchange effects. Both a fixed-rate loan at 5\% and a floating-rate loan
whose rate happened to be 5\% generate the same entries:
\begin{center}
\begin{tabular}{lrr}
\toprule Account & Debit & Credit\\\midrule
Cash, on initial borrowing &100&\\
Borrowing liability &&100\\\midrule
Interest expense, at year-end &5&\\
Cash &&5\\\bottomrule
\end{tabular}
\end{center}
Even the full historical journal would not distinguish these two contracts.
Now impose a scenario in which the floating rate resets immediately to 7\%
for the whole next year. The fixed loan produces interest of 5; the floating
loan produces interest of 7. If the disclosure is also compatible with any
floating-rate fraction $\alpha\in[0,1]$, using the convention of
Chapter~\ref{sim:observations}, future interest is
\begin{equation}\label{eq:running}
 I(\alpha)=100[0.05(1-\alpha)+0.07\alpha]=5+2\alpha.
\end{equation}
The historical accounts are compatible with the whole interval $[5,7]$ of
future interest. This is a range of feasible amounts. It becomes a probability
distribution only when a distribution for $\alpha$ is supplied. Neither the
number of generated interpretations nor a language model's confidence provides
that distribution.

Journals remain essential because simulated economic events must have coherent
accounting consequences. However, the simulated state must also contain
contract terms, maturities, customer or borrower cohorts, collateral,
unrecognised commitments, and management decisions. These determine what the
journals will be under a changed economy. Reconstructing every historical
transaction is unnecessary when a smaller state preserves the responses needed
for the intended scenarios. Establishing that smaller state is an explicit
modelling task, examined in Chapter~\ref{ch:math}.

The analogy with competing legal interpretations is useful at the level of
representation. Several readings may survive the evidence and require separate
reasoning. In accounting, some alternatives concern the meaning of a disclosure;
others concern undisclosed facts or future behaviour. The system must identify
which kind of uncertainty is involved. A known contractual term must not become
an arbitrary interpretive choice merely because another choice is convenient
for a simulation.

\section{Users, questions and deliverables}

An accounting analyst needs to inspect a number in its original table, verify
its entity, currency and period, and see how it reconciles with the notes. A
financial modeller needs to distinguish a reported balance from assumptions
about its composition. A language engineer needs deterministic semantics and
tests. A researcher needs to compare explanations under the same scenario,
including explanations that fit history equally well.

The first release must answer the following complete workflow. An analyst
loads a set of reports and specifies an information cutoff. AccDSL preserves
the original documents and extracts source-linked observations. It identifies
reconciliation failures and constructs an explicitly bounded family of
compatible histories. The analyst declares additional assumptions where
simulation requires them. The system executes each requested scenario and
returns statements, journal traces, explanation-specific distributions or
bounds, and unresolved limitations. When a later filing arrives, the analyst
can compare the revised interpretation with the earlier information set.

\begin{longtable}{L{34mm}L{102mm}}
\toprule Requirement & Required behaviour\\\midrule\endhead
\req{R1 Evidence} & Every observed amount or narrative assertion retains its
document version, location, extraction method and original representation.\\
\req{R2 History} & Queries distinguish economic time, reporting period,
publication time and database ingestion time; restatements remain recoverable.\\
\req{R3 Meaning} & Entity, perimeter, dimensions, currency, measurement basis,
standard version and accounting concept accompany values.\\
\req{R4 Ambiguity} & Discrete alternatives and continuous unknowns retain their
joint constraints across accounts and periods. Missing values never mean zero.\\
\req{R5 Accounting} & All committed journals balance in the book's functional
currency, and balances reconcile to explicitly classified movements.\\
\req{R6 Economics} & Contracts, exposures and behavioural assumptions drive
events. Accounting identities alone never determine scenario responses.\\
\req{R7 Results} & Conditional results survive by interpretation and scenario.
Weighted results require declared, normalised probability measures.\\
\req{R8 Explanation} & A report cell can be traced to journals, events, rules,
assumptions and evidence, including consolidation adjustments.\\
\req{R9 Extensibility} & Standard, industry and company packages extend shared
types through versioned interfaces, with explicit unsupported cases.\\
\req{R10 Reproduction} & A saved run records source versions, model versions,
parameters, calendars, seeds and numerical conventions.\\\bottomrule
\end{longtable}

These are acceptance requirements. Python, SQL storage, a particular parser,
and a particular constraint solver are implementation choices. Initial support
will cover a non-financial company with cash, receivables, payables, inventory,
fixed assets, ordinary debt and ordinary equity. Banks and insurers must be
representable in the core data model from the outset, but their executable
measurement and behavioural packages follow dedicated milestones. Until a
package supports an instrument, the result is explicitly incomplete.

\section{What constitutes a useful first implementation}

The first demonstrator should reconstruct three reporting dates and the
intervening annual and interim flows for one company, including at least one
restatement or revised comparative. It must preserve a borrowing-composition
ambiguity and execute the fixed/floating example through both journals and all
affected statement lines. A synthetic company supplies known hidden history;
a public company supplies realistic disclosure and extraction problems. Their
roles differ: the synthetic case tests preservation and rejection of known
alternatives, while the public case tests traceability and practical coverage.

Success is a reproducible chain of evidence and computation, with identified
gaps. A convincing-looking complete ledger reconstructed from incomplete
disclosures would fail this product contract if its assumptions were hidden.
Forecast accuracy requires a later evaluation with historical information
cutoffs and suitable baselines; it is not established by the first demonstrator.
